% paper/sections/10b_spatial.tex
%
% §10.1  空間離散化の精度検証
%   CCD/DCCD 収束性・散逸特性・曲率精度・ψ直接法・ε_eff 診断
% ═══════════════════════════════════════════════════════════════════════════════
\subsection{空間離散化の精度検証}
\label{sec:verify_spatial}
% ═══════════════════════════════════════════════════════════════════════════════

% ─────────────────────────────────────────────────────────────────────────────
\subsubsection{CCD/DCCD 微分演算の収束性}
\label{sec:verify_ccd_convergence}
% ─────────────────────────────────────────────────────────────────────────────

CCD 法（§\ref{sec:ccd_def}）の $\Ord{h^6}$ 精度を定量的に確認する．
周期境界条件下でテスト関数 $f = \sin(2\pi x)\sin(2\pi y)$ を用い，
$N = 16, 32, 64, 128, 256$ の各格子解像度で 1 階・2 階微分の $L_2$/$L_\infty$ 誤差を計測する．
1 階微分の $L_2$ 誤差は $N = 16 \to 128$ の範囲で精確に $\Ord{h^6}$ の収束を示す
（収束次数 $p \approx 6.00$--$6.06$）．
2 階微分も同等の $\Ord{h^{5.97}}$ 精度を達成する．
$N = 256$ では打ち切り誤差が $\Ord{10^{-14}}$ に到達し，
丸め誤差のフロアに接触するため見かけ上の収束次数が低下する
（$L_\infty$ で $p \approx 5.36$）．
これは CCD の理論限界ではなく倍精度浮動小数点の限界である．

\begin{table}[htbp]
\centering
\caption{CCD 1 階微分の格子収束（周期境界，$f = \sin(2\pi x)\sin(2\pi y)$）}
\label{tab:ccd_periodic}
\begin{tabular}{rrrrr}
\toprule
$N$ & $L_2$ 誤差 & 次数 & $L_\infty$ 誤差 & 次数 \\
\midrule
 16 & $1.28 \times 10^{-6}$  & ---  & $2.57 \times 10^{-6}$  & ---  \\
 32 & $1.93 \times 10^{-8}$  & 6.06 & $3.86 \times 10^{-8}$  & 6.06 \\
 64 & $2.99 \times 10^{-10}$ & 6.01 & $5.97 \times 10^{-10}$ & 6.01 \\
128 & $4.65 \times 10^{-12}$ & 6.00 & $9.36 \times 10^{-12}$ & 6.00 \\
256 & $7.67 \times 10^{-14}$ & 5.92 & $2.28 \times 10^{-13}$ & 5.36 \\
\bottomrule
\end{tabular}
\end{table}

% ─────────────────────────────────────────────────────────────────────────────
\subsubsection{DCCD の数値散逸特性}
\label{sec:verify_dccd_dissipation}
% ─────────────────────────────────────────────────────────────────────────────

DCCD（§\ref{sec:dccd_decoupling}）の散逸メカニズムを
離散伝達関数 $H(\xi; \varepsilon_d)$ により検証する．
伝達関数は $e^{i\xi x}$ に対する CCD の応答比として定義され，
DCCD の場合（§\ref{sec:dccd_decoupling}）
\begin{equation}
  H(\xi; \varepsilon_d) = 1 - 4\varepsilon_d \sin^2(\xi/2)
  \label{eq:dccd_transfer_verify}
\end{equation}
と解析的に表される．Nyquist 波数 $\xi = \pi$ での減衰率を表~\ref{tab:dccd_transfer} に示す．

\begin{table}[htbp]
\centering
\caption{DCCD 伝達関数の Nyquist 波数における値 $H(\pi; \varepsilon_d)$}
\label{tab:dccd_transfer}
\begin{tabular}{rr}
\toprule
$\varepsilon_d$ & $H(\pi)$ \\
\midrule
0.00 & $1.0000$（散逸なし：CCD）\\
0.05 & $0.8000$（移流用：低散逸）\\
0.25 & $0.0000$（PPE 用：完全抑制）\\
\bottomrule
\end{tabular}
\end{table}

$\varepsilon_d = 1/4$ のとき $H(\pi) = 0$ となり，
Nyquist モード（チェッカーボード振動）が\textbf{完全に抑制}される．
これは §\ref{sec:checkerboard_solution} の理論的結論と一致する．
$\varepsilon_d = 0.05$（移流ステップ用）では $H(\pi) = 0.80$ であり，
高波数成分を 20\% 減衰させつつ，低波数の移流精度を保持する．

数値実験として $N = 128$ 格子上でチェッカーボードモード $(-1)^{i+j}$ に
DCCD フィルタ（$\varepsilon_d = 0.25$）を適用した結果，
フィルタ後の振幅はゼロとなることを確認した．

% ─────────────────────────────────────────────────────────────────────────────
\subsubsection{\texorpdfstring{曲率 $\kappa$ の計算精度}{曲率 κ の計算精度}}
\label{sec:verify_curvature}
% ─────────────────────────────────────────────────────────────────────────────

曲率計算（§\ref{sec:curvature}，式~\eqref{eq:curvature_2d}）の精度を
2 種類のテスト問題と 3 つの計算経路で検証する．
3 経路の定義を以下に示す：
\begin{itemize}
  \item \textbf{経路 A（$\phi$ 直接）：}
        真の符号付距離関数 $\phi$ に CCD を適用し $\kappa$ を計算（理想ベースライン）．
  \item \textbf{経路 B（$\psi \to \phi$ ロジット逆変換）：}
        CLS 変数 $\psi$ からロジット逆変換
        $\phi = \varepsilon\ln[\psi/(1-\psi)]$（§\ref{sec:newton_inversion}）で
        $\phi$ を復元し，$\phi$ の CCD 微分から $\kappa$ を計算（従来法）．
  \item \textbf{経路 C（$\psi$ 直接法，§\ref{sec:curvature_direct_psi_cls}）：}
        $\psi$ の CCD 微分値から式~\eqref{eq:curvature_psi_2d} で $\kappa$ を直接計算
        （ロジット逆変換を省略；§\ref{sec:curvature_direct_psi_cls} の推奨経路）．
\end{itemize}

\textbf{(a) 正弦波状界面 $y = 0.5 + A\sin(2\pi x)$，$A = 0.05$：}
曲率の解析解 $\kappa = f''/(1+f'^2)^{3/2}$ と CCD 計算値を比較する．
経路 A では $N = 32 \to 256$ で $L_\infty$ 誤差の収束次数は
$p \approx 4.96$--$5.00$ を示す（表~\ref{tab:curvature_sinusoidal}）．
これは CCD 6 次精度微分から期待される高次収束と整合する．

\begin{table}[htbp]
\centering
\caption{正弦波状界面の曲率 $L_\infty$ 誤差収束（$A = 0.05$，$\varepsilon = 1.5h$）}
\label{tab:curvature_sinusoidal}
\begin{tabular}{rrrrrrr}
\toprule
$N$ & A: $\phi$ 直接 & 次数 & B: ロジット経由 & 次数 & C: $\psi$ 直接 & 次数 \\
\midrule
 32 & $7.64 \times 10^{-4}$ & ---  & $7.64 \times 10^{-4}$ & ---  & $3.84 \times 10^{-2}$ & ---  \\
 64 & $2.45 \times 10^{-5}$ & 4.96 & $2.14 \times 10^{-5}$ & 5.16 & $2.19 \times 10^{-2}$ & 0.81 \\
128 & $7.72 \times 10^{-7}$ & 4.99 & $5.80 \times 10^{-6}$ & 1.88 & $2.27 \times 10^{-2}$ & $-0.05$ \\
256 & $2.42 \times 10^{-8}$ & 5.00 & $1.33 \times 10^{-5}$ & $-1.20$ & $3.59 \times 10^{-2}$ & $-0.66$ \\
\bottomrule
\end{tabular}
\end{table}

経路 A は全格子で $\Ord{h^5}$ を維持する．
経路 B は $N \leq 64$ で A と同等の精度を示すが，$N \geq 128$ では
ロジット逆変換の飽和域処理に起因する誤差が顕在化して収束が停滞し，
$N = 256$ では誤差が増大する（収束次数が負値）．
経路 C は全格子で収束しない（$\varepsilon = \Ord{h}$ でのシグモイド勾配の桁落ちによる）．

円形界面（$R = 0.25$，$\kappa_\mathrm{exact} = 4$）では，
経路 B の曲率収束は CCD・CD2 ともに $\Ord{h^1}$ にとどまり，
ロジット逆変換の界面厚み $\varepsilon = \Ord{h}$ が律速する．
CCD 微分演算自体は (a) で $\Ord{h^5}$ を達成しており，精度の上界ではない．

本稿では経路 B（$\psi \to \phi$ ロジット逆変換→CCD）を採用する．
経路 C（$\psi$ 直接法）は $\varepsilon = \Ord{h}$ でシグモイド勾配の桁落ちにより
格子精緻化で収束しないため不採用とした（表~\ref{tab:curvature_sinusoidal}）．
なお，GFM 界面処理（§\ref{sec:gfm}）の精度が $\Ord{h^2}$（界面近傍）で律速されるため，
曲率計算経路の選択は全体精度に影響しない．

% ─────────────────────────────────────────────────────────────────────────────
\subsubsection{非一様格子上の微分精度と幾何学的保存則（GCL）}
\label{sec:gcl_verification}
% ─────────────────────────────────────────────────────────────────────────────
% paper/sections/11b4_gcl_verification.tex
%
% §11.3.4  非一様格子上の微分精度と幾何学的保存則（GCL）検証
%          実験スクリプト: experiments/gcl_verification.py

\paragraph{主張．}

本実験が示す主命題は以下の二点である：

\begin{center}
\textbf{%
  界面適合格子（$\alpha>1$）上でも CCD の漸近収束次数は一様格子と同等であり，%
  定数速度場（フリーストリーム）に対する偽の空間微分は計算機精度に抑えられる．}
\end{center}

第~\ref{sec:grid}章の界面適合格子は，計算座標 $\xi=i/N$（一様）を Jacobian
$J_i=\partial\xi/\partial x$ を介して物理座標 $x$ に写像する．
CCD は均一な $\xi$-空間で解を構成し，以下の連鎖則（§~\ref{sec:CCD}）で物理微分に変換する：
\begin{equation}
  \frac{\partial f}{\partial x} = J\,\frac{\partial f}{\partial\xi}, \qquad
  \frac{\partial^2 f}{\partial x^2}
    = J^2\,\frac{\partial^2 f}{\partial\xi^2}
      + J\,\frac{\mathrm{d}J}{\mathrm{d}\xi}\,\frac{\partial f}{\partial\xi}.
  \label{eq:gcl_chain}
\end{equation}

\paragraph{実験設定．}

\textbf{Test 1（収束次数の保全）．}
$f(x)=\sin(\pi x)$（解析微分既知）を格子点で評価し，$N=16,32,64,128$ で格子を精緻化する．
一様格子（$\alpha=1$）と非一様格子（$\alpha=2$，界面を $x=0.5$ に配置）の
$L_\infty$ 誤差を比較する．

\textbf{Test 2（GCL / フリーストリーム保存）．}
定数場 $f=1$ を与え，CCD が生成する偽の空間微分 $\|\partial(1)/\partial x\|_\infty$
を計算する．ナビエ・ストークス方程式の散逸項 $\nabla\cdot\bm{u}$ は一次微分演算子であるため，
一次誤差が計算機精度 $\varepsilon_\mathrm{mach}$ に収まることを合格基準とする：
\begin{equation}
  \left\|\frac{\partial(1)}{\partial x}\right\|_\infty \leq 10^3\,\varepsilon_\mathrm{mach}.
  \label{eq:gcl_criterion}
\end{equation}

\paragraph{結果．}

\begin{table}[htbp]
\centering
\caption{非一様格子（$\alpha=2$）と一様格子上の CCD 微分精度：
  $f=\sin(\pi x)$，$L_\infty$ 誤差と収束次数．}
\label{tab:gcl_convergence}
\begin{tabular}{c
                r@{\ }l r@{\ }l
                r@{\ }l r@{\ }l}
\toprule
& \multicolumn{4}{c}{一様格子 ($\alpha=1$)}
& \multicolumn{4}{c}{非一様格子 ($\alpha=2$)} \\
\cmidrule(lr){2-5}\cmidrule(lr){6-9}
$N$
& \multicolumn{2}{c}{$E_{d1}$} & \multicolumn{2}{c}{次数}
& \multicolumn{2}{c}{$E_{d1}$} & \multicolumn{2}{c}{次数} \\
\midrule
 16 & $6.0\!\times\!10^{-5}$  & & $-$   & & $4.3\!\times\!10^{-2}$  & & $-$   & \\
 32 & $9.6\!\times\!10^{-7}$  & & $5.96$& & $7.4\!\times\!10^{-4}$  & & $5.86$& \\
 64 & $1.5\!\times\!10^{-8}$  & & $5.99$& & $7.3\!\times\!10^{-6}$  & & $6.65$& \\
128 & $2.4\!\times\!10^{-10}$ & & $6.00$& & $1.1\!\times\!10^{-7}$  & & $6.08$& \\
\bottomrule
\end{tabular}
\end{table}

\begin{tcolorbox}[resultbox, title={Test 1 結果：収束次数の保全}]
非一様格子（$\alpha=2$）上でも $d/dx$ の漸近収束次数は $O(h^6)$ に収束し，
一様格子と同等の次数を示す（表~\ref{tab:gcl_convergence}）．
$N=16$ における絶対誤差の差は，格子圧縮領域での前漸近的効果によるものであり，
$N \geq 32$ では一様格子と同様の高次収束が得られる．
\end{tcolorbox}

\begin{tcolorbox}[resultbox, title={Test 2 結果：GCL 合格}]
非一様格子（$\alpha=2$，$N=16\text{--}128$）上で $f=1$ を与えた場合，
一次偽微分は式~\eqref{eq:gcl_criterion} の基準を満たす：
\[
  \max_{N}\left\|\frac{\partial(1)}{\partial x}\right\|_\infty
    = 5.4\!\times\!10^{-14}
    \ll 10^3\,\varepsilon_\mathrm{mach} = 2.2\!\times\!10^{-13}.
\]
二次微分 $\partial^2(1)/\partial x^2$ は Jacobian 項 $J\,(\mathrm{d}J/\mathrm{d}\xi)\,\varepsilon_\mathrm{mach}$
に起因する~$O(N\,\varepsilon_\mathrm{mach})$ の残差を持つが，
これはナビエ・ストークス方程式の散逸項（一次演算子）とは無関係であり，
GCL の合否に影響しない．
\end{tcolorbox}


% ─────────────────────────────────────────────────────────────────────────────
\subsubsection{DCCD フィルタ：1次元移流問題による散逸特性の定量評価}
\label{sec:dccd_dissipation_bench}
% ─────────────────────────────────────────────────────────────────────────────
% paper/sections/11b6_dccd_dissipation.tex
%
% §11.3.6  DCCD フィルタ：分散・散逸特性と選択的高周波減衰の評価
%          実験スクリプト: experiments/dccd_comparison.py
%          結果データ:    results/dccd_comparison/

\paragraph{主張．}

本実験が示す主命題は以下の二点である：

\begin{center}
\textbf{%
  DCCD の 10 次選択フィルタ（§~\ref{sec:advection_dccd_design}）は，%
  不連続プロファイルに対して CCD の過大な変動（全変動：TV）を大幅に抑制しつつ，%
  滑らかなプロファイルでは CCD と同等の高精度を維持する（選択的高周波散逸）．%
}
\end{center}

CLS 法では界面 $\psi$ が $0$–$1$ 遷移層を持つ準不連続プロファイルである．
CCD（フィルタなし）は高波数成分に対して分散誤差が蓄積し，
界面近傍に非物理的な振動（Gibbs 型）を生じやすい．
DCCD フィルタはスペクトル領域で高波数成分のみを選択的に減衰させるため，
滑らかな領域の精度を犠牲にせずに振動を抑制する（§~\ref{sec:advection_dccd_design}）．

\paragraph{実験設定．}

1次元線形移流方程式
\begin{equation}
  \frac{\partial u}{\partial t} + c\,\frac{\partial u}{\partial x} = 0,
  \qquad x \in [0,1),\quad c = 1,\quad T = 1
  \label{eq:linear_advection_dccd}
\end{equation}
を $N=256$ の周期格子，CFL $=0.4$，RK4 で 1 周期（$T=1$）移流する．
初期条件として 3 種類のプロファイルを用いる：
\begin{itemize}[noitemsep]
  \item \textbf{Square}：不連続矩形パルス（$x \in [0.3, 0.5)$ で $u=1$，他は $0$）
  \item \textbf{Triangle}：$C^0$ 三角波（頂点 $x=0.4$ で $u=1$，$[0.3,0.5]$ 台）
  \item \textbf{Smooth tanh}：双曲線正接型平滑化矩形（準不連続・滑らか）
\end{itemize}
比較スキーム：O2（2次中心差分），O4（4次中心差分），CCD（6次），
DCCD（CCD $+$ 10次選択フィルタ，$\alpha_f=0.4$），WENO5（5次重み付き ENO）．

評価指標：
$L_2 = \sqrt{\text{mean}(u-u_\text{exact})^2}$，
$\text{TV} = \sum_i |u_{i+1}-u_i|$（周期ラップ込み）．

\paragraph{結果．}

\begin{table}[htbp]
\centering
\caption{1周期後の $L_2$ 誤差（$N=256$，CFL$=0.4$，太字：各初期条件での最良値）．
  スクリプト: \texttt{experiments/dccd\_comparison.py}．}
\label{tab:dccd_l2_bench}
\begin{tabular}{lrrrrr}
\toprule
初期条件 & O2 & O4 & CCD & DCCD & WENO5 \\
\midrule
Square（不連続） & $1.41\times10^{-1}$ & $9.08\times10^{-2}$ & $4.80\times10^{-2}$ & $\mathbf{4.23\times10^{-2}}$ & $6.35\times10^{-2}$ \\
Triangle（$C^0$） & $2.41\times10^{-2}$ & $5.84\times10^{-3}$ & $\mathbf{1.37\times10^{-3}}$ & $1.41\times10^{-3}$ & $3.90\times10^{-3}$ \\
Smooth tanh & $3.75\times10^{-2}$ & $2.46\times10^{-3}$ & $\mathbf{2.57\times10^{-5}}$ & $3.75\times10^{-5}$ & $9.66\times10^{-4}$ \\
\bottomrule
\end{tabular}
\end{table}

\begin{table}[htbp]
\centering
\caption{1周期後の全変動（TV）．
  Exact 値は初期条件の TV（Square: 2.000, Triangle: 1.969, tanh: 2.000）．
  太字：各初期条件での最良値（TV $\to$ exact に最近）．}
\label{tab:dccd_tv_bench}
\begin{tabular}{lrrrrr}
\toprule
初期条件 & O2 & O4 & CCD & DCCD & WENO5 \\
\midrule
Square（不連続） & $18.40$ & $17.17$ & $10.83$ & $3.15$ & $\mathbf{2.00}$ \\
Triangle（$C^0$） & $2.60$ & $2.37$ & $2.09$ & $1.97$ & $\mathbf{1.91}$ \\
Smooth tanh & $2.82$ & $2.02$ & $2.00$ & $2.00$ & $\mathbf{2.00}$ \\
\bottomrule
\end{tabular}
\end{table}

\begin{figure}[htbp]
  \centering
  \includegraphics[width=0.95\linewidth]{%
    ../results/dccd_comparison/figures/04_waveforms_panel.png}
  \caption{%
    3 種の初期条件で $T=1$ 後の波形比較（O2・O4・CCD・DCCD・WENO5）．
    左から Square, Triangle, tanh．
    Square では CCD に顕著な Gibbs 型振動が残るが，DCCD フィルタにより大幅に抑制される．
    DCCD の TV は $3.15$（CCD の $10.83$ に対して約 1/3）であり，WENO5 の $2.00$ に近い．
    Smooth tanh では CCD と DCCD の波形はほぼ重なり，フィルタが滑らかな成分に影響しないことを示す
    （$L_2$: CCD $2.57\times10^{-5}$ vs DCCD $3.75\times10^{-5}$）．}
  \label{fig:dccd_waveforms_bench}
\end{figure}

\begin{tcolorbox}[resultbox, title={DCCD フィルタの選択的散逸効果}]
\begin{itemize}[noitemsep]
  \item \textbf{不連続（Square）：} DCCD は CCD に比べ TV を $10.83 \to 3.15$ に低減（約 $1/3$）し，
        Gibbs 型振動を大幅に抑制する．$L_2$ 誤差も最良（$4.23\times10^{-2}$）．
        WENO5 が TVD 条件を満たすため最良 TV（$2.00$）だが，$L_2$ は DCCD より悪い（$6.35\times10^{-2}$）．
  \item \textbf{平滑（Smooth tanh）：} DCCD と CCD の $L_2$ 誤差は同オーダー
        （$3.75\times10^{-5}$ vs $2.57\times10^{-5}$，1.5 倍以内），TV も同値（$2.000$）．
        フィルタが高波数成分のみを対象とし，滑らかな成分を保護することを示す．
  \item \textbf{CLS 法への含意：} CLS の界面 $\psi$ は Smooth tanh に相当する平滑化プロファイルであるため，
        DCCD は CCD と同等の精度で界面を移流しつつ，界面近傍のわずかな非物理的振動を抑制する．
        これは §~\ref{sec:interface_advection_bench}（Zalesak・単一渦）の低体積誤差の理論的根拠となる．
\end{itemize}
\end{tcolorbox}


